\begin{recipe}{Saucijzenbroodjes}
\kicker[Bakken,Vlees,Nederlands,Snack]{Bakken · vlees · Nederlands · snack}
\index[register]{Saucijzenbroodjes}
\index[register]{Bladerdeeg}
\index[register]{Half-om-half gehakt}

\meta{45 minuten \dvd 6 stuks \dvd Oven 200\,°C}

\lettrine{S}{aucijzenbroodjes} passen net zo goed op een verjaardag als op een
gewone avond. Met een eigen kruidenmengsel smaken ze merkbaar beter dan de
versie uit de diepvries, en zes stuks zijn zo klaar.

\blockrule{Ingrediënten}
\begin{ingredients}
  \ing{250 g}{half-om-half gehakt}
  \ing{30 g}{paneermeel}
  \ing{6 plakjes}{bladerdeeg}
  \ing{1}{ei}
  \ing{5 g}{gehaktkruiden (zie hieronder)}
\end{ingredients}

\medskip
\noindent\textit{Gehaktkruiden (maakt ca.\ 5\,g):}
\begin{ingredients}
  \ing{2 g}{zout}
  \ing{1 g}{paprikapoeder (mild of gerookt)}
  \ing{0{,}7 g}{zwarte peper, gemalen}
  \ing{0{,}6 g}{uienpoeder}
  \ing{0{,}4 g}{knoflookpoeder}
  \ing{0{,}3 g}{nootmuskaat}
  \ing{0{,}2 g}{korianderpoeder}
\end{ingredients}

\blockrule{Bereiding}
\tip{Bladerdeeg werkt het best als het koud blijft. Laat het niet langer ontdooien dan nodig.}
\begin{steps}
  \item Laat het gehakt en de bladerdeegplakjes ontdooien.
  \item Verwarm de oven voor op 200\,°C (boven- en onderwarmte).
  \item Splits het ei: tik de schaal in tweeën en schenk het eigeel heen en
        weer tussen de twee schaalhelften boven een kommetje, tot al het
        eiwit in het kommetje is gelopen en alleen het eigeel achterblijft in
        de schaal. Houd een beetje eiwit apart voor het bestrijken; klop de
        rest los en voeg dit samen met de gehaktkruiden en het paneermeel bij het
        gehakt.
  \item Kneed het vleesmengsel goed door tot alles gelijkmatig verdeeld is.
  \item Weeg porties van 60\,g gehakt af en leg elke portie in een rol op de
        helft van een plakje bladerdeeg.
  \item Bestrijk de vrije rand van het deeg dun met het achtergehouden ei en
        vouw het deeg dubbel over het vlees, zodat de natte rand op de
        andere rand plakt. Druk de naad goed dicht.
  \item Dip een vork kort in water, zodat hij niet aan het deeg blijft
        plakken, en druk er de kenmerkende streepjes mee in de gesloten
        rand.
  \item Leg de broodjes met de naad naar onderen op een met bakpapier beklede
        bakplaat en bestrijk de bovenkant met het achtergehouden ei.
  \item Bak 30 minuten op 200\,°C tot de broodjes goudbruin zijn.
\end{steps}

\heroimagefade{images/sauzijcenbroodjes}
\end{recipe}
